% MQO: a reta de mínimos quadrados de um AR(1) no log-retorno da PETR4
% (o caminho q = 0 do CArima, que é uma regressão só) e a soma dos quadrados
% dos resíduos para cada coeficiente, com o mínimo no estimado.
\documentclass[border=4pt]{standalone}
% Preâmbulo comum dos gráficos que explicam os métodos (30/10/2026).
% Dados e números: dados/resultados/metodos_ilustrados_2026-10-30/, gerados por
% ferramentas/MetodosIlustrados.mq5 com as funções do próprio robô. Nenhum
% número é digitado aqui: os do texto vêm de valores.tex (\vX para o texto,
% \pX para coordenada).
\usepackage[T1]{fontenc}
\usepackage{amsmath}
\usepackage{mathptmx}  % Times, como as figuras do TG
\usepackage{pgfplots}
\usepgfplotslibrary{groupplots,dateplot,colormaps,fillbetween}
\usetikzlibrary{arrows.meta,positioning}
\pgfplotsset{compat=1.18}
\definecolor{cmod}{RGB}{0,90,181}     % modelo / ajuste
\definecolor{cdes}{RGB}{220,50,32}    % destaque, limite, valor crítico
\definecolor{cdad}{RGB}{120,120,120}  % dado observado
\newcommand{\arq}[1]{../../dados/resultados/metodos_ilustrados_2026-10-30/#1}
\input{\arq{valores.tex}}
% Figuras 08 a 14 (o robô operando): traço por barra da refação de 05/10/2026,
% resumido por ferramentas/scripts/robo_em_acao.sh. Macros \vR e \pR.
\newcommand{\arqr}[1]{../../dados/resultados/robo_em_acao_2026-10-30/#1}
\input{\arqr{valores.tex}}
\definecolor{cari}{RGB}{230,140,0}   % ARIMA puro
\definecolor{cok}{RGB}{0,140,70}     % compra liberada
\usetikzlibrary{shapes.misc}
\newcommand{\passo}[1]{\tikz[baseline=(n.base)]\node[circle,fill=black!75,text=white,inner sep=1pt,font=\footnotesize\bfseries](n){#1};}
% eixo de datas comum às figuras do traço (jul/2019 a jun/2021)
\pgfplotsset{traco/.style={date coordinates in=x, xmin=2019-07-01, xmax=2021-06-30,
    xtick={2019-07-01,2020-01-01,2020-07-01,2021-01-01,2021-06-30},
    xticklabels={jul/2019,jan/2020,jul/2020,jan/2021,jun/2021}}}
\pgfplotsset{
  base/.style={table/col sep=comma, grid=major, grid style={gray!20},
    tick label style={font=\small}, label style={font=\small},
    title style={font=\small\bfseries, align=center},
    every axis plot/.append style={line width=0.9pt},
    /pgf/number format/use comma, /pgf/number format/1000 sep={.}},
}
% faixa vertical nos pregões do traço em que a condição #1 vale, na cor #2,
% do piso #3 ao teto #4 do eixo
\newcommand{\faixa}[4]{%
  \addplot[name path=topo, draw=none, const plot, forget plot] table[x=data, y expr={(#1) ? #4 : #3}]{\arqr{dia_a_dia.csv}};
  \addplot[name path=piso, draw=none, forget plot] coordinates {(2019-07-01,#3) (2021-06-30,#3)};
  \addplot[#2, forget plot] fill between[of=topo and piso];}

\begin{document}
\begin{tikzpicture}
\begin{groupplot}[base, group style={group size=2 by 1, horizontal sep=1.9cm},
  width=7.6cm, height=6.4cm]
\nextgroupplot[title={Cada ponto é um pregão:\\retorno de ontem $\times$ retorno de hoje},
  xlabel={Retorno de ontem, $r_{t-1}$ (\%)}, ylabel={Retorno de hoje, $r_t$ (\%)},
  xmin=-8, xmax=8, ymin=-8, ymax=8]
\addplot[cdad!45, line width=0.3pt, quiver={u=0, v=\thisrow{ajuste}-\thisrow{y}, every arrow/.append style={-}}]
  table[x=x, y=y]{\arq{mqo_dispersao.csv}};
\addplot[only marks, mark=*, mark size=0.9pt, cdad, fill opacity=0.6, draw opacity=0.6]
  table[x=x, y=y]{\arq{mqo_dispersao.csv}};
\addplot[cmod, line width=1.6pt, domain=-8:8, samples=2] {\pMqoC + \pMqoPhi*x};
\node[anchor=north west, font=\small, fill=white, fill opacity=0.85, text opacity=1, align=left]
  at (rel axis cs:0.02,0.98)
  {$\hat r_t = \vMqoC + \vMqoPhi\, r_{t-1}$\\[1pt]
   {\color{cdad}traços: os resíduos $r_t - \hat r_t$}};
\nextgroupplot[title={O MQO escolhe o coeficiente que\\minimiza a soma dos quadrados},
  xlabel={Coeficiente $\phi$ testado}, ylabel={Soma dos quadrados dos resíduos (\%$^2$)},
  xmin=-0.5, xmax=0.5, xtick={-0.4,-0.2,0,0.2,0.4},
  /pgf/number format/fixed]
\addplot[cmod, line width=1.4pt] table[x=phi, y=sqr]{\arq{mqo_perfil.csv}};
\draw[cdes, dashed, line width=0.9pt] (axis cs:\pMqoPhi,\pgfkeysvalueof{/pgfplots/ymin}) -- (axis cs:\pMqoPhi,\pMqoSQR);
\addplot[only marks, mark=*, mark size=2.6pt, cdes] coordinates {(\pMqoPhi,\pMqoSQR)};
\node[anchor=south, font=\small, align=center, text=cdes] at (axis cs:\pMqoPhi,\pMqoSQR)
  {mínimo\\$\hat\phi = \vMqoPhi$};
\end{groupplot}
\end{tikzpicture}
\end{document}
